\chapter{Sprint 65 --- Cụm J Phần V: Việt hoá source, V0 $\to$ ★VR (2026-10-05 $\to$ 10-07)}

\emph{Chapter này ghi lại Phần V của cụm J: đổi toàn bộ câu gốc tiếng Việt viết cứng trong source code team sang tiếng Anh,
còn bản tiếng Việt chuyển vào \texttt{i18n/vi\_VN.po}. Người dùng tiếng Việt vẫn thấy đúng chữ cũ. Người dùng tiếng Anh, Thái,
Trung không còn thấy tiếng Việt lẫn vào giao diện. Có 10 phiên code (V1 \ldots\ V8b) và một phiên review (★VR). Nhật ký ở
\texttt{docs/f-progress.md}; quy ước, bài học từng phiên ở \texttt{docs/next-session-clusters-J.md} §6; quy trình công cụ ở
\texttt{scripts/qa/README.md} §Phần V.}

\section{Vì sao có Phần V}

Ngày 05/10, khi làm tool dịch \texttt{wujia\_i18n} (J-T1+T2), máy quét \texttt{scripts/qa/vn\_hardcode\_scan.py} đếm được
\textbf{2\,659 chuỗi tiếng Việt viết cứng}: code team khoảng 2\,375, code anh Thái 284, còn lại 945 chuỗi trong test.
Odoo chỉ dịch được khi câu gốc là tiếng Anh. Vì vậy dù đã có tool dịch, portal ở en/th/zh vẫn hiện tiếng Việt.
Chủ dự án chốt làm Phần V \textbf{trước} J-T4 (``ưu tiên xử lý trước'').

\section{Quy ước chốt ở J-V0}

\begin{enumerate}[nosep]
\item Câu gốc trong QWeb/Python/JS đổi sang tiếng Anh ngắn, giữ nghĩa nghiệp vụ. Thêm một dòng vào \texttt{docs/i18n-glossary.csv}
  (key = câu EN, VN = câu cũ), rồi \texttt{sync\_translations.py} sinh \texttt{vi\_VN.po} + \texttt{.pot}.
\item Hằng nhãn Python dùng \texttt{\_lt()}, dịch lúc render theo env người xem. Thông báo controller dùng \texttt{\_()}.
\item JS portal không có \texttt{\_t}: câu đặt trong \texttt{<t t-set>} (dịch được), gắn lên phần tử bằng
  \texttt{data-wj-msg-*}, JS đọc qua \texttt{wjMsg(el, key, fallbackEN)}.
\item User tiếng Anh thấy tiếng Anh, user tiếng Việt thấy tiếng Việt. Khách chưa đăng nhập trên \texttt{/portal} mặc định vi\_VN.
  Bật đủ vi/en/th/zh\_CN.
\item BA không duyệt câu EN từng phiên. Gom một danh sách gửi BA một lần ở ★VR. Bản zh/th sẽ dịch máy ở J-T5.
\item Mỗi phiên phải đạt: quét lại module = 0; \texttt{wj\_text\_probe} vi\_VN 0 lệch chữ; en/th không vỡ layout; suite 0 đỏ mới;
  so bản dịch trong DB 0 chỗ VN $\to$ EN.
\end{enumerate}

\section{Các phiên}

\begin{center}
\footnotesize
\begin{tabularx}{\textwidth}{>{\raggedright\arraybackslash}p{1.2cm}>{\raggedright\arraybackslash}p{1.6cm}>{\raggedright\arraybackslash}X>{\raggedright\arraybackslash}p{1.3cm}}
\toprule
Phiên & Commit & Module (số chuỗi) & Ngày \\
\midrule
V0 & \texttt{d2966119} & Quy ước + công cụ (\texttt{vn\_to\_en\_pairs}, \texttt{wj\_text\_probe}, mốc vi\_VN 52 trang) & 05/10 \\
V1 & \texttt{7b03b680} & \texttt{portal\_layout} (259) + helper \texttt{wjMsg} & 05/10 \\
V2 & \texttt{7e1d5abf} & \texttt{portal\_base} (382); màu badge theo câu EN & 05/10 \\
V3 & \texttt{baa5b547} & \texttt{portal\_exam} (345) + \texttt{exam} (63) & 06/10 \\
V4 & \texttt{359fe36c}, \texttt{6129ca59} & \texttt{portal\_sale} (201) + \texttt{sale} (16) + \texttt{order\_window} (8) & 06/10 \\
V5 & \texttt{2f9f5695} & \texttt{portal\_debt} (196) + \texttt{account} (4) & 06/10 \\
V6 & \texttt{b78c4568} & \texttt{portal\_return} (168) + \texttt{return} (69), migration seed loại lỗi & 06/10 \\
V7 & \texttt{351932a8} & support (112+2), knowledge (46+2), notification (86+27); WJ-SUPPORT-003 & 07/10 \\
V8a & \texttt{dbd4335b} & \texttt{purchase\_history} (99), \texttt{portal\_delivery} (93), \texttt{delivery} (15), \texttt{fleet} (11) & 07/10 \\
V8b & \texttt{02e363fc} & \texttt{info\_request} (85+7), \texttt{portal\_report} (69), \texttt{core} (9), metabase (1) & 07/10 \\
★VR & chưa commit & Review cả Phần V (mục dưới) & 07/10 \\
\bottomrule
\end{tabularx}
\end{center}

V3 \ldots\ V8b đã lên UAT (lần cuối 07/10, \texttt{6ef4db78}); chủ dự án duyệt nâng tay \texttt{wujia\_metabase\_connector}.

\section{★VR --- review cả Phần V}

Đo trên DB \texttt{wujia\_vr}, chép từ \texttt{wujia\_t1} (mốc trước V1, cùng dữ liệu với \texttt{docs/i18n-baseline/vi\_VN.json}),
rồi \texttt{-u} 26 module team.

\begin{center}
\footnotesize
\begin{tabularx}{\textwidth}{>{\raggedright\arraybackslash}p{4.2cm}>{\raggedright\arraybackslash}X}
\toprule
Thước & Kết quả \\
\midrule
Quét chuỗi VN viết cứng & Code team \textbf{0} trên 24 module; anh Thái 284 (gửi danh sách) \\
Chữ không dấu máy quét sót & 1 chỗ: cột ``SL'' ở chi tiết chuyến giao mobile $\to$ ``Qty'' (vi giữ ``SL'') \\
msgid vi\_VN mới chưa dịch & 83, không chỗ nào người dùng thấy (gallery dev, template chết \texttt{signup}/\texttt{forgot\_pass\_back}, field backend, metabase) \\
\texttt{wj\_text\_probe} vi\_VN, 54 trang & 2 lần đo 0 lệch nhau; so mốc trước Phần V: chỉ thêm \texttt{data-wj-msg-*} + 2 sửa có chủ đích (tiêu đề khung giờ thi V4, danh mục hỗ trợ V7) $\Rightarrow$ cập nhật mốc \\
DB vi\_VN \texttt{wujia\_t1} $\leftrightarrow$ \texttt{wujia\_vr} & field/help/selection/model/menu/action/arch: \textbf{0 chỗ VN $\to$ EN}, arch không mất term tiếng Việt \\
Playwright vi/en/th $\times$ 1440/390 & 156 trang: 0 tràn ngang, 0 lỗi JS, 0 HTTP $\geq$ 400; chữ Việt còn trên en/th = dữ liệu + ``Tiếng Việt'' + 2 menu Khảo sát (module Thái) \\
Độ phủ vi\_VN (\texttt{wujia\_i18n}) & Module portal 96,6--100\,\% (thiếu = chữ giống nhau 2 ngôn ngữ: PDF, ID, Email, chip ``i''); backend L2 41--95\,\% do field backend vốn tiếng Anh từ trước \\
Suite 25 module (trừ core) & 1\,185 test: 2 đỏ $\to$ 1 sửa (test bảng BA), 1 do DB đo có sẵn kết quả quét cũ (\texttt{wujia\_i18n}); nhóm liên quan chạy lại 544/0 \\
Mutation & 4/4 bị bắt \\
\bottomrule
\end{tabularx}
\end{center}

\subsection{Bỏ nhánh tra ngược nhãn tiếng Việt}

Từ V2, \texttt{status\_badge\_for(label)} có một nhánh tạm. Nếu nhãn là tiếng Việt, nhánh này tra ngược \texttt{vi\_VN.po} ra câu EN
để tô màu badge, cho các module chưa Việt hoá. Rà 8 module gọi hàm này thì thấy 2 chỗ vẫn truyền nhãn \emph{đã dịch}: Lịch sử đặt
hàng (danh sách + chi tiết) và màn ``Đặt hàng thành công'' mobile. Ở vi\_VN, 2 chỗ này đúng màu chỉ nhờ nhánh tạm. Với th/zh sau
J-T5, nhãn sẽ là tiếng Thái/Trung và badge rơi về màu xám.

Sửa: Lịch sử lấy màu qua \texttt{portal\_order\_badge(order)} sẵn có. Màn kết quả dùng helper mới
\texttt{portal\_order\_state\_badge(state)}, tính màu từ \texttt{\_lt}. Sau đó xoá nhánh tra ngược khỏi code chạy thật. Test chống lùi
cần ``màu của nhãn tiếng Việt trước phiên'' làm đáp án, nên logic tra ngược chuyển sang \texttt{tests/common.py::legacy\_vn\_badge}
(chỉ test dùng). Version: \texttt{portal\_base} 19.0.7.37.1, \texttt{portal\_sale} 19.0.5.1.1, \texttt{purchase\_history} 19.0.4.0.1,
\texttt{portal\_delivery} 19.0.4.0.1.

\subsection{Danh sách gửi đi}

\begin{itemize}[nosep]
\item \texttt{docs/i18n-review/ba-en-terms.csv}: 1\,361 cặp câu VN cũ $\to$ EN mới cho BA rà một lần. 10 dòng cần chốt: một câu
  tiếng Việt đang có 2 câu tiếng Anh, ví dụ ``Khẩn cấp'' là \emph{Emergency} ở loại thông báo nhưng là \emph{Critical} ở ticket.
\item \texttt{docs/i18n-review/thai-vn-hardcode.csv}: 284 chuỗi trong module anh Thái (18 dòng có gợi ý EN từ glossary).
\item Báo anh Thái: \texttt{\#: web\_survey\_ui} trong \texttt{.po} của \texttt{wujia\_franchise\_inspection} làm mỗi lần \texttt{-u}
  ghi 120 dòng ERROR; WJ-INSPECT-001 (trang Khảo sát không theo ngôn ngữ phiên).
\end{itemize}

\section{Bài học của cả Phần V}

\begin{itemize}[nosep]
\item Hàm tô màu theo \emph{nhãn} phải nhận câu gốc (\texttt{\_lt}), không nhận chữ đã dịch. Nếu không, màu chỉ đúng ở ngôn ngữ
  có bảng tra ngược. Trả \texttt{(nhãn đã dịch, màu)} cùng một chỗ, đừng để caller tự tô từ nhãn đã dịch.
\item Máy quét nhận tiếng Việt qua dấu, nên sót chữ không dấu và viết tắt (\texttt{Xem}, \texttt{Trang}, \texttt{SL}, \texttt{SP},
  \texttt{T2}\ldots\texttt{CN}). Sau mỗi \texttt{sync} phải liệt kê msgid chưa dịch và grep danh sách viết tắt quen.
\item Glossary dùng chung ghi đè bản dịch riêng của module. Luôn so \texttt{.po} cũ/mới bằng babel (chỉ đọc) và so bản dịch trong DB
  (0 chỗ VN $\to$ EN), sửa \texttt{.po} bằng text.
\item Bản dịch \texttt{model:} gắn theo xmlid: DB cũ vẫn hiện chữ Việt dù \texttt{.po} mới bỏ trống. Chỉ phép so DB mới bắt được.
\item Đáp án test chống lùi phải là bảng viết cứng \emph{trước} phiên. So bằng \texttt{env.\_} cùng ngôn ngữ thì đột biến ``bỏ
  \texttt{load\_vi}'' vẫn xanh.
\item File xuất (XLSX), câu JS, \texttt{aria-label}, tiêu đề tab đều là chữ người dùng thấy, phải đi đường dịch như QWeb.
\end{itemize}

\section{Nợ để lại}

\begin{itemize}[nosep]
\item Deploy ★VR: \texttt{-u wujia\_portal\_base,wujia\_portal\_sale,wujia\_portal\_purchase\_history,wujia\_portal\_delivery} + restart.
\item Backend L2 (support, knowledge, fleet, info\_request, exam, notification\ldots) còn field/view tiếng Anh chưa có bản vi (có từ
  trước Phần V, không phải lùi). Làm bằng J-T5 hoặc một lượt riêng nếu chủ dự án cần backend tiếng Việt đủ.
\item Template chết \texttt{wujia\_portal\_layout.signup\_form} / \texttt{forgot\_pass\_back} và trang gallery dev vẫn nằm trong
  \texttt{.pot}. Xoá khi dọn layout.
\item 284 chuỗi của anh Thái; zh/th thấy EN tới J-T5.
\end{itemize}
